% TEMPLATE-MILC-v3.tex - plantilla maestra para todas las guias MILC v3
% Ver scripts/generadores/build-guia-milc.py para la lista completa de
% claves esperadas. El script valida que no queden placeholders sin
% reemplazar antes de escribir el archivo final.

\documentclass[11pt,letterpaper]{article}
\usepackage[margin=1.75cm]{geometry}
\usepackage[table]{xcolor}
\usepackage{fontspec}
\usepackage[spanish]{babel}
\usepackage{tikz}
\usepackage[most]{tcolorbox}
\usepackage{tabularx}
\usepackage{array}
\usepackage{fancyhdr}
\usepackage{titlesec}
\usepackage{ragged2e}
\usepackage{enumitem}
\usepackage{microtype}

\setmainfont{Helvetica Neue}
\setsansfont{Helvetica Neue}
\setlength{\parindent}{0pt}
\setlength{\parskip}{6pt}
\hyphenpenalty=200
\exhyphenpenalty=200
\pretolerance=400
\tolerance=2000
\setlength{\emergencystretch}{1em}
\renewcommand{\arraystretch}{1.25}
\setlist{leftmargin=5mm,itemsep=1mm,topsep=1mm}
\newcolumntype{Y}{>{\RaggedRight\arraybackslash}X}

\definecolor{milcMagenta}{HTML}{E6007E}
\definecolor{milcVino}{HTML}{5A0038}
\definecolor{milcTurquesa}{HTML}{0093A5}
\definecolor{milcVerde}{HTML}{58A923}
\definecolor{milcMostaza}{HTML}{E5B400}
\definecolor{milcOcre}{HTML}{B86600}
\definecolor{milcNegro}{HTML}{241816}
\definecolor{milcGris}{HTML}{F7F7F4}
\definecolor{milcLinea}{HTML}{E8E2DD}
\definecolor{milcGradePrimary}{HTML}{FF2D87}
\definecolor{milcGradeDark}{HTML}{9F1B56}
\definecolor{milcGradeSoft}{HTML}{FFE0EE}
\definecolor{milcGradeAccent}{HTML}{CFE7FF}
\definecolor{milcGradeText}{HTML}{FFFFFF}
\color{milcNegro}

\pagestyle{fancy}
\fancyhf{}
\lhead{\small Tecnología e Informática · Grado 10}
\rhead{\small Guía 20 · MILC v3}
\cfoot{\small\thepage}
\renewcommand{\headrulewidth}{0pt}

\newtcolorbox{titlebox}[1]{
  enhanced,colback=#1,colframe=#1,arc=7mm,boxrule=0pt,
  left=7mm,right=7mm,top=3mm,bottom=3mm,
  before skip=8pt,after skip=8pt,
  fontupper=\sffamily\bfseries\Large\color{white}
}

\newtcolorbox{softbox}[3]{
  enhanced,colback=#2,colframe=#1,borderline west={4pt}{0pt}{#1},
  arc=2mm,boxrule=.4pt,left=4mm,right=4mm,top=3mm,bottom=3mm,
  before skip=6pt,after skip=8pt,title={#3},coltitle=white,
  colbacktitle=#1,fonttitle=\sffamily\bfseries,fontupper=\RaggedRight,
  attach boxed title to top left={xshift=3mm,yshift=-2mm},
  boxed title style={arc=2mm, boxrule=0pt}
}

\newtcolorbox{drawbox}[2]{
  enhanced,colback=white,colframe=#1,arc=4mm,boxrule=1pt,
  left=5mm,right=5mm,top=4mm,bottom=4mm,before skip=8pt,after skip=8pt,
  title={#2},coltitle=white,colbacktitle=#1,fonttitle=\sffamily\bfseries,
  fontupper=\RaggedRight,
  attach boxed title to top left={xshift=4mm,yshift=-2mm},
  boxed title style={arc=3mm, boxrule=0pt}
}

\newtcolorbox{infoband}[1]{
  enhanced,colback=#1!8,colframe=#1!50,arc=2mm,boxrule=.5pt,
  left=4mm,right=4mm,top=2mm,bottom=2mm,before skip=4pt,after skip=4pt,
  fontupper=\small\RaggedRight
}

% Puente narrativo entre secciones (contrato editorial Conéctate).
% Da continuidad: "Acabas de X. Ahora vas a Y porque Z."
% Visualmente: banda izquierda en color del grado, texto en itálica pequeña.
\newtcolorbox{puentebox}{
  enhanced,colback=milcGradeSoft,colframe=milcGradeDark,
  borderline west={3pt}{0pt}{milcGradeDark},
  arc=0mm,boxrule=0pt,left=5mm,right=4mm,top=2.5mm,bottom=2.5mm,
  before skip=4pt,after skip=8pt,
  fontupper=\itshape\small\color{milcNegro}\RaggedRight
}
\newcommand{\puente}[1]{%
  \begin{puentebox}%
    \textcolor{milcGradeDark}{\textbf{\textrightarrow}}\hspace{2mm}#1%
  \end{puentebox}%
}

\newcommand{\linead}[1]{\noindent\color{milcLinea}\rule{#1}{0.5pt}\color{milcNegro}}
\newcommand{\respuesta}[1]{\par\vspace{2mm}\foreach \n in {1,...,#1}{\noindent\color{milcLinea}\rule{\linewidth}{0.5pt}\par\vspace{5mm}}\color{milcNegro}}
\newcommand{\checkbox}{\textcolor{milcGradeDark}{\fbox{\phantom{x}}}\hspace{5pt}}

\newcommand{\phasecard}[4]{
\begin{tcolorbox}[enhanced,colback=#3,colframe=#2,arc=3mm,boxrule=.5pt,height=3.05cm,valign=center,halign=center,left=1.5mm,right=1.5mm,top=2mm,bottom=2mm]
{\sffamily\bfseries\fontsize{22}{24}\selectfont\textcolor{#2}{#1}}\par
{\sffamily\bfseries\small #4}
\end{tcolorbox}
}

\begin{document}
\thispagestyle{empty}
\begin{tikzpicture}[remember picture,overlay]
  \fill[white] (current page.south west) rectangle (current page.north east);
  \path[fill=milcGradePrimary,rounded corners=16mm]
    ([xshift=.55cm,yshift=-.55cm]current page.north west) rectangle
    ([xshift=-.55cm,yshift=3.65cm]current page.south east);

  \node[anchor=north west,text=milcGradeText] at ([xshift=2.0cm,yshift=-1.85cm]current page.north west)
    {\sffamily\bfseries\fontsize{14}{16}\selectfont GRADO};
  \node[anchor=north west,text=milcGradeText,opacity=.82] at ([xshift=2.0cm,yshift=-2.75cm]current page.north west)
    {\sffamily\bfseries\fontsize{9}{11}\selectfont SERIE GUÍAS · TECNOLOGÍA E INFORMÁTICA};

  \begin{scope}[shift={([xshift=-3.05cm,yshift=-2.55cm]current page.north east)}]
    \fill[white,opacity=.80] (-.62,-.52) rectangle (.62,.52);
    \draw[milcGradeText,line width=1pt] (-.62,-.52) rectangle (.62,.52);
    \fill[milcGradeDark] (-.42,-.38) rectangle (-.22,.06);
    \fill[milcGradeAccent] (-.10,-.38) rectangle (.10,.24);
    \fill[milcGradePrimary!60!black] (.22,-.38) rectangle (.42,.38);
  \end{scope}

  \node[anchor=west,text=milcGradeText] at ([xshift=1.9cm,yshift=-12.85cm]current page.north west)
    {\sffamily\bfseries\fontsize{124}{124}\selectfont 10};
  \draw[milcGradeText,line width=1.7pt]
    ([xshift=4.52cm,yshift=-11.68cm]current page.north west) circle (.13cm);

  \node[anchor=west,text=milcGradeText,text width=8.7cm,align=left] at ([xshift=10.35cm,yshift=-11.6cm]current page.north west)
    {
     {\sffamily\bfseries\fontsize{19}{22}\selectfont Sustentación P2 ---\\defensa pública\\del oficio aprendido}\\[4mm]
     {\sffamily\bfseries\fontsize{23}{26}\selectfont Guía 20-10-TIC}\\[2mm]
     {\sffamily\fontsize{13}{16}\selectfont Período 2 · Informes técnicos y comunicación profesional con IA}\\[5mm]
     {\sffamily\bfseries\fontsize{11}{13}\selectfont Institución Educativa Sor María Juliana}\\[1mm]
     {\sffamily\fontsize{10}{12}\selectfont Autor: Dr. Álvaro Cárdenas Orozco}
    };

  \path[fill=milcGradeSoft,rounded corners=7mm]
    ([xshift=1.35cm,yshift=.85cm]current page.south west) rectangle
    ([xshift=-1.35cm,yshift=4.25cm]current page.south east);
  \draw[milcGradeDark,line width=.65pt,rounded corners=7mm]
    ([xshift=1.35cm,yshift=.85cm]current page.south west) rectangle
    ([xshift=-1.35cm,yshift=4.25cm]current page.south east);
  \node[anchor=center,text=milcNegro,text width=17.0cm,align=center] at ([yshift=2.55cm]current page.south)
    {\sffamily\fontsize{10}{12}\selectfont Metodología MILC · Apertura ancestral · Pensamiento computacional · Triángulo Dussel-Estoicismo-Floridi\\
    \textcolor{milcGradeDark}{\bfseries Producto:} Sustentación pública de 5 minutos del mini-proyecto integrador del periodo + 4-5 slides simples como apoyo + demo de fragmento clave del informe técnico + sesión de Q\&A de 5 minutos con honestidad ante preguntas técnicas + autoevaluación final con 2 fortalezas concretas y 2 mejoras específicas para el periodo siguiente};
\end{tikzpicture}
\mbox{}
\newpage

\section*{Apertura: Sustentación del periodo 2 --- defensa del oficio}

\begin{softbox}{milcGradeDark}{milcGradeSoft}{Saber ancestral}
En los gremios tradicionales del oficio (carpinteros, ebanistas, sastres, alfareros), hubo durante siglos un acto ritual que cualquier aprendiz aceptaba antes de recibir título de oficial: \textbf{la defensa pública de la pieza maestra ante los mayores del gremio}. El aprendiz no entregaba la pieza en silencio: la \textbf{defendía}. Se sentaba ante el maestro y los oficiales viejos, traía su pieza, la presentaba, explicaba cómo la había hecho, qué decisiones tomó, qué aprendió. Los mayores hacían preguntas técnicas y a veces incómodas: \emph{``¿por qué ese ensamble así?''}, \emph{``¿qué pasa si la madera trabaja?''}, \emph{``¿cómo responderías a un cliente que reclama por defecto a los 6 meses?''}. El aprendiz no podía decir \emph{``no sé''} sin perder la prueba: debía defender cada decisión con argumento técnico, ético, económico. Si la defensa era sólida, los mayores le entregaban el título de oficial. Si flaqueaba en algún punto, volvía al taller un año más. La sabiduría era inquebrantable: \textbf{el oficio se demuestra hablando, no solo haciendo}. Esa práctica ancestral del gremio es exactamente lo que la sustentación pública del periodo 2 te pide hoy.
\end{softbox}

\begin{softbox}{milcTurquesa}{milcGris}{Saber contemporáneo}
La \textbf{sustentación pública del periodo} es el acto formal donde el estudiante defiende el mini-proyecto integrador (sesión 9) ante el grupo. Tiene 3 propiedades irrenunciables: (1) \textbf{Tiempo controlado}: 5 minutos (similar a sustentación TED corta). Obliga a destilar lo esencial del periodo entero. (2) \textbf{Estructura clara}: tema del mini-proyecto + proceso + hallazgos clave + uso de IA declarado + recomendaciones. (3) \textbf{Honestidad ante Q\&A}: cuando la audiencia hace preguntas duras, responder con verdad técnica. La estructura típica de los 5 minutos: (a) \textbf{Introducción del proyecto} (1 min): qué es y para quién. (b) \textbf{Datos y metodología} (1 min): encuesta aplicada con cuántas respuestas. (c) \textbf{Hallazgos clave} (1 min): los 2-3 resultados más importantes con cifras. (d) \textbf{Recomendaciones} (1 min): qué propones y a quién. (e) \textbf{Declaración honesta de IA} (1 min): qué modelos, porcentajes, qué hiciste tú. Después: 5 minutos de Q\&A. La regla profesional: \textbf{lo que no puedes contar en 5 minutos, todavía no lo dominas del todo}.
\end{softbox}

\begin{softbox}{milcMostaza}{milcGris}{Pregunta-puente}
\textit{¿Qué sabía el aprendiz del gremio al sentarse a defender su pieza ante los mayores, que el novato olvida cuando confunde sustentar con leer slides? ¿Y por qué la declaración honesta de uso de IA en la sustentación es más profesional que pretender autoría pura?}
\end{softbox}

\begin{softbox}{milcVerde}{milcGris}{Saber y saber hacer}
Al terminar podrás: (1) \textbf{analizar} tu mini-proyecto identificando los 2-3 puntos más fuertes para destacar y los 1-2 puntos débiles que aún sostienen revisión; (2) \textbf{explicar} en sustentación pública la concepción, el proceso, los hallazgos y el uso de IA; (3) \textbf{evaluar} con honestidad preguntas duras del grupo sin inventar respuestas; (4) \textbf{crear} las 4-5 slides simples y la autoevaluación final con compromisos para el periodo 3.
\end{softbox}

\small
\begin{tabularx}{\textwidth}{>{\bfseries}p{3.0cm}Y}
\rowcolor{milcGradePrimary!12}Autor & Dr. Álvaro Cárdenas Orozco\\
DBA & Comunicación profesional con asistencia de IA y herramientas ofimáticas gratuitas (MEN, T\&I)\\
Referentes & Markdown como estándar abierto · Google Workspace · Estoicismo (claridad)\\
Producto final & Sustentación pública de 5 minutos del mini-proyecto integrador del periodo + 4-5 slides simples como apoyo + demo de fragmento clave del informe técnico + sesión de Q\&A de 5 minutos con honestidad ante preguntas técnicas + autoevaluación final con 2 fortalezas concretas y 2 mejoras específicas para el periodo siguiente\\
\end{tabularx}
\normalsize

\puente{Antes de empezar, mira el viaje completo. Has recorrido 9 sesiones del periodo 2. Hoy es el \textbf{día de defender lo aprendido}. Después de hoy, abre el periodo 3 (ofimática con IA, contabilidad, emprendimiento) que será el puente directo al grado 11.}

\begin{titlebox}{milcGradeDark}
Ruta MILC: así vamos a aprender
\end{titlebox}
\begin{tabularx}{\textwidth}{>{\centering\arraybackslash}X>{\centering\arraybackslash}X>{\centering\arraybackslash}X>{\centering\arraybackslash}X}
\phasecard{1}{milcVerde}{milcGris}{Escuta\\[-1mm]\small Observo, escucho y pregunto.} &
\phasecard{2}{milcTurquesa}{milcGris}{Sistematización\\[-1mm]\small Pensamiento computacional.} &
\phasecard{3}{milcMagenta}{milcGris}{Praxis\\[-1mm]\small Construyo y aplico.} &
\phasecard{4}{milcVino}{milcGris}{Evaluación\\[-1mm]\small Reflexión liberadora.}
\end{tabularx}


\puente{Antes de teorizar, vas a hacer un \emph{inventario honesto} de tu mini-proyecto: 2-3 puntos fuertes (lo que estás orgulloso de haber logrado) y 1-2 puntos débiles (lo que aún necesita revisión). Esa lista cruda es la base de la sustentación.}

\begin{titlebox}{milcVerde}
Fase 1 · Escuta
\end{titlebox}

\begin{softbox}{milcVerde}{milcGris}{Escena para escuchar}
\textbf{Actividad 1 · ANALIZA --- Inventario honesto del mini-proyecto} (10 min · individual). Toma tu mini-proyecto de la sesión 9 y anota con sinceridad: (a) \textbf{2-3 puntos fuertes}: lo que estás orgulloso de haber logrado (datos sólidos, análisis bien auditado, propuesta clara, voz propia en la carta del autor). (b) \textbf{1-2 puntos débiles}: lo que aún necesita revisión (muestra pequeña, alguna recomendación sin responsable, IA con porcentaje alto que no se declaró bien). Esa lista cruda es la materia prima honesta de la sustentación.
\end{softbox}

\checkbox Anoté 2-3 puntos fuertes específicos.\\
\checkbox Anoté 1-2 puntos débiles honestos.\\
\checkbox Reconozco la importancia de declarar ambos.

\begin{infoband}{milcVerde}
\textbf{Lo que va al cuaderno (Actividad 1 · ANALIZA).} Título: «Actividad 1 --- Inventario honesto». \textbf{Formato:} 2 listas (fuertes, débiles). \textbf{Sabes que terminaste cuando} podrías declarar ambas en sustentación pública.
\end{infoband}

\puente{Ya tienes inventario. Ahora vas a entender la \textbf{anatomía de la sustentación de 5 minutos + Q\&A}: estructura, las 5 preguntas duras más típicas, los 5 errores del sustentante novato.}

\begin{titlebox}{milcTurquesa}
Fase 2 · Sistematización con pensamiento computacional
\end{titlebox}

\begin{softbox}{milcTurquesa}{milcGris}{Cuatro pilares aplicados al tema}
Tu inventario es la base. Ahora necesitas la \textbf{anatomía profesional} de la sustentación de 5 minutos + Q\&A.
\end{softbox}

\small
\begin{tabularx}{\textwidth}{>{\bfseries\RaggedRight\arraybackslash}p{3.6cm}Y}
\rowcolor{milcTurquesa!90}\color{white}Pilar & \color{white}\bfseries Aplicación al tema\\
1. Descomposición & La \textbf{estructura de la sustentación de 5 minutos} se descompone así (1 minuto por bloque): (1) \textbf{Introducción del proyecto}: nombre del mini-proyecto, tema, destinatario al que se dirige. (2) \textbf{Metodología y datos}: encuesta aplicada, cuántas respuestas, herramientas usadas. (3) \textbf{Hallazgos clave}: 2-3 hallazgos numéricos importantes con interpretación. (4) \textbf{Recomendaciones}: qué propones, a quién, con qué plazo. (5) \textbf{Declaración honesta de IA}: modelos usados, porcentajes aproximados, qué hiciste tú. Después: \textbf{Q\&A de 5 minutos} con 3-5 preguntas del grupo.\\
2. Reconocimiento de patrones & Las \textbf{5 preguntas duras} más típicas que recibirás: (a) \emph{``¿cuántas respuestas tuvo tu encuesta?''} (responder con cifra). (b) \emph{``¿cómo verificaste que la IA no se equivocó al clasificar?''} (responder con tabla auditoría). (c) \emph{``¿tu recomendación es realmente factible para el destinatario?''} (defender con argumentos). (d) \emph{``¿qué porcentaje del informe escribió la IA?''} (responder con honestidad numérica). (e) \emph{``¿qué cambiarías si tuvieras otra semana?''} (responder con autoevaluación honesta).\\
3. Abstracción & Los \textbf{5 errores típicos} del sustentante novato: (a) \textbf{Leer las slides}: matar la sustentación en el primer minuto. (b) \textbf{Ocultar uso de IA}: pretender autoría pura cuando IA fue significativa. (c) \textbf{Inflar hallazgos}: presentar 18 respuestas como \emph{``estudio amplio''}. (d) \textbf{Evadir preguntas duras}: \emph{``buena pregunta...''} y desviar el tema. (e) \textbf{Cierre apurado}: terminar con \emph{``gracias''} en lugar de invitar al destinatario a recibir el informe.\\
4. Algoritmo conceptual & Algoritmo: (1) preparar 4-5 slides simples (1 por bloque del minuto). (2) escribir guion hablado de 700-800 palabras (= 5 min). (3) ensayar 2 veces cronometrado. (4) sustentar. (5) responder Q\&A con honestidad. (6) autoevaluar con 2 fortalezas + 2 mejoras + 1 compromiso para periodo 3.\\
\end{tabularx}
\normalsize

\begin{drawbox}{milcTurquesa}{Anatomía de la sustentación P2}
\textbf{1 min:} Introducción del proyecto. \\[1mm] \textbf{1 min:} Metodología y datos. \\[1mm] \textbf{1 min:} Hallazgos clave con cifras. \\[1mm] \textbf{1 min:} Recomendaciones. \\[1mm] \textbf{1 min:} Declaración honesta de IA. \\[1mm] \textbf{5 min Q\&A:} honestidad ante preguntas duras.
\end{drawbox}

\begin{softbox}{milcMostaza}{milcGris}{Errores comunes y su corrección}
(1) Leer las slides. (2) Ocultar uso de IA. (3) Inflar hallazgos. (4) Evadir preguntas duras. (5) Cierre apurado. (6) Pasarse del tiempo. (7) Slides saturadas.
\end{softbox}

\begin{infoband}{milcTurquesa}
\textbf{Lo que va al cuaderno (Actividad 2 · EXPLICA).} Título: «Actividad 2 --- Anatomía sustentación». \textbf{Formato:} plan minuto a minuto + 5 preguntas duras con respuesta + 7 errores. \textbf{Sabes que terminaste cuando} podrías sustentar mañana.
\end{infoband}

\puente{Tienes el método. Toca aplicarlo: armas slides simples, ensayas cronometrado, sustentas, respondes Q\&A, autoevalúas con compromisos para periodo 3.}

\begin{titlebox}{milcMagenta}
Fase 3 · Praxis con pensamiento computacional
\end{titlebox}

\begin{softbox}{milcMagenta}{milcGris}{Construyendo el producto con los cuatro pilares}
\textbf{Actividad 3 · EVALÚA + CREA --- Sustentación ejecutada + autoevaluación} (45 min · individual + colectivo). Hora de ejecutar. Preparas slides, ensayas, sustentas, respondes Q\&A, autoevalúas con compromisos.
\end{softbox}

\small
\begin{tabularx}{\textwidth}{>{\bfseries\RaggedRight\arraybackslash}p{3.6cm}Y}
\rowcolor{milcMagenta!90}\color{white}Pilar & \color{white}\bfseries Aplicación al producto\\
Descomposición & Tu sustentación se construye en 6 pasos: (1) preparar 4-5 slides simples (1 por bloque del minuto, máximo 25 palabras por slide). (2) escribir guion hablado de 700-800 palabras. (3) ensayar cronometrado 2 veces. (4) sustentar con honestidad técnica. (5) responder Q\&A con verdad. (6) autoevaluar con 2 fortalezas + 2 mejoras + 1 compromiso para el periodo 3 (ofimática y emprendimiento).\\
Patrones & La \textbf{declaración honesta de uso de IA} es el momento crucial de la sustentación. Modelo: \emph{``Usé ChatGPT y Gemini para borradores iniciales del informe (capítulos 2-4 del mini-proyecto, aproximadamente 50 por ciento del texto generado), edité a mano todas las conclusiones y la carta del autor (50 por ciento restante), audité la clasificación de la IA respuesta por respuesta corrigiendo 6 de 18, y declaré modelos en la bitácora''}. Esa precisión gana respeto profesional.\\
Abstracción & La \textbf{autoevaluación final} tiene 5 elementos: (a) 2 fortalezas concretas (manejo del tiempo, claridad de los hallazgos, honestidad en declaración de IA). (b) 2 mejoras específicas (más ensayo, mejor manejo del nervio, slides más simples). (c) 1 compromiso concreto para el periodo 3: \emph{``En el periodo 3 voy a llegar con más práctica de presentación oral''}, \emph{``En el periodo 3 voy a aplicar el análisis con IA desde el primer dato, no al final''}.\\
Algoritmo de armado & Algoritmo: (1) preparar slides. (2) ensayar. (3) sustentar. (4) responder. (5) autoevaluar. (6) entregar al docente: archivo slides + guion + autoevaluación.\\
\end{tabularx}
\normalsize

\begin{drawbox}{milcMagenta}{Checklist antes de sustentar}
\checkbox 4-5 slides simples (máx. 25 palabras cada una). \\[1mm] \checkbox Guion hablado de 700-800 palabras. \\[1mm] \checkbox 2 ensayos cronometrados hechos. \\[1mm] \checkbox Declaración honesta de IA preparada. \\[1mm] \checkbox 5 preguntas duras con respuesta preparada. \\[1mm] \checkbox Sustentación dentro de 5 minutos. \\[1mm] \checkbox Autoevaluación con 2+2+1 elementos.
\end{drawbox}

\begin{softbox}{milcMagenta}{milcGris}{Plantilla de trabajo}
\textbf{1 min Intro.} \textit{Mini-proyecto + destinatario.} \\[1mm] \textbf{1 min Metodología.} \textit{Encuesta + N respuestas.} \\[1mm] \textbf{1 min Hallazgos.} \textit{2-3 con cifras.} \\[1mm] \textbf{1 min Recomendaciones.} \textit{Qué + a quién + plazo.} \\[1mm] \textbf{1 min IA.} \textit{Modelos + porcentajes + lo humano.} \\[1mm] \textbf{Q\&A 5 min.} \textit{Honestidad ante 5 preguntas duras.}
\end{softbox}

\begin{infoband}{milcMagenta}
\textbf{Lo que va al cuaderno (Actividad 3 · EVALÚA + CREA).} Título: «Actividad 3 --- Mi sustentación P2». \textbf{Formato:} (a) archivo de slides; (b) guion hablado; (c) preguntas recibidas con respuestas; (d) autoevaluación con compromiso. \textbf{Sabes que terminaste cuando} entregaste materiales y reservaste tiempo para reflexión.
\end{infoband}

\puente{Lo que sigue resume \emph{qué} entregas hoy: sustentación + slides + autoevaluación. El criterio principal: que la audiencia entienda lo que aprendiste y qué propones, en 5 minutos.}

\begin{titlebox}{milcMostaza}
Producto de la sesión
\end{titlebox}
\textbf{Sustentación 5 min + 4-5 slides + Q\&A + autoevaluación}.
\begin{softbox}{milcMostaza}{milcGris}{Criterios de calidad}
(1) Estructura de 5 minutos respetada. (2) 4-5 slides simples. (3) Declaración honesta de IA con porcentajes. (4) Q\&A respondida con verdad técnica. (5) Autoevaluación con 2 fortalezas + 2 mejoras + 1 compromiso. (6) Compromiso específico para periodo 3.
\end{softbox}

\puente{Antes de cerrar, mira la sustentación desde las cinco dimensiones humanas. Defender el oficio aprendido con honestidad es respeto por las 9 sesiones del periodo; pasar la sustentación sin asumir lo que falta es traicionar al estudiante que fuiste.}

\begin{titlebox}{milcVino}
Fase 4 · Evaluación liberadora — cinco dimensiones
\end{titlebox}

\begin{softbox}{milcVino}{milcGris}{Sentido de la evaluación}
Evaluar no es solo revisar una nota. Es reconocer qué aprendí, qué cambió en mí, qué emociones manejé, cómo me relaciono con la ciudad y la comunidad, y qué le dejaré a quien viene detrás.
\end{softbox}

\textbf{Mi reflexión liberadora en cinco dimensiones.} Completa cada frase iniciada:

\small
\begin{tabularx}{\textwidth}{>{\bfseries\RaggedRight\arraybackslash}p{3.4cm}Y>{\RaggedRight\arraybackslash}p{4.6cm}}
\rowcolor{milcVino}\color{white}Dimensión & \color{white}\bfseries Mi respuesta (frame starter) & \color{white}\bfseries Algo que descubrí\\
1. Personal & Lo que más cambió en mí fue \linead{6.5cm} & \linead{4.2cm}\\
2. Emocional & La emoción más fuerte fue \linead{2.5cm} y la regulé \linead{3cm} & \linead{4.2cm}\\
3. Ciudadana & En democracia esto importa porque \linead{6cm} & \linead{4.2cm}\\
4. Local (Cartago) & En mi barrio sería útil para \linead{6.5cm} & \linead{4.2cm}\\
5. Intergeneracional & A mi abuela/abuelo le diría \linead{6.5cm} & \linead{4.2cm}\\
\end{tabularx}
\normalsize

\begin{infoband}{milcVino}
\textbf{Para la próxima sustentación.} Vuelve a esta tabla en seis meses. Verás que las respuestas cambian — no porque lo aprendido cambie, sino porque tú cambias. La evaluación liberadora es una conversación contigo a través del tiempo.
\end{infoband}


\puente{Y para terminar, tres voces sobre el cierre del periodo. Dussel sobre el aprendiz que devuelve oficio a la comunidad, Epicteto sobre la honestidad ante preguntas duras, Floridi sobre la sustentación con declaración de IA como nueva ética profesional.}

\begin{titlebox}{milcGradeDark}
Triángulo de pensamiento · cierre filosófico
\end{titlebox}

\begin{softbox}{milcVino}{milcGris}{Dussel · lente del nosotros}
\textit{``El aprendiz que devuelve oficio asumido a la comunidad cierra el ciclo de la pedagogía liberadora.''}

Tu sustentación cierra el periodo 2: lo aprendido vuelve al grupo hecho voz propia. La filosofía de la liberación pide ese acto. Cuando declaras con honestidad lo que hiciste y lo que aún te falta, estás devolviendo oficio asumido. La phronesis del cierre consiste en \textbf{compartir el aprendizaje, no guardarlo}.

\textbf{¿Mi sustentación devuelve oficio al grupo, o solo cumple con requisito?}
\end{softbox}
\linead{\linewidth}\\[1mm]
\linead{\linewidth}

\begin{softbox}{milcMostaza}{milcGris}{Estoicismo · Epicteto · cuidado interior}
\textit{``La honestidad ante preguntas duras es virtud profesional; evadir es debilidad disfrazada de cortesía.''}

Cuando alguien te pregunta \emph{``¿qué porcentaje del informe lo escribió la IA?''}, es tentador minimizar. La disciplina estoica recomienda lo opuesto: \emph{declara los porcentajes exactos con sobriedad}. Esa precisión gana respeto que el evasivo nunca gana. La phronesis del cierre honra la verdad numérica.

\textbf{¿Estoy respondiendo preguntas duras con honestidad numérica, o evadiendo por timidez?}
\end{softbox}
\linead{\linewidth}\\[1mm]
\linead{\linewidth}

\begin{softbox}{milcTurquesa}{milcGris}{Floridi · lente de la infoesfera}
\textit{``La sustentación con declaración honesta de IA es el modelo profesional del siglo XXI.''}

En la era de IA generativa, las sustentaciones profesionales que declaran abiertamente uso de IA con porcentajes y modelos se posicionan como ético emergente. Tu sustentación de hoy te entrena en esa práctica profesional que la industria adoptará obligatoriamente en pocos años.

\textbf{¿Mi sustentación contribuye al modelo transparente del siglo XXI, o sostiene la opacidad heredada?}
\end{softbox}
\linead{\linewidth}\\[1mm]
\linead{\linewidth}

\begin{titlebox}{milcVino}
Compromiso final
\end{titlebox}

\small
\begin{tabularx}{\textwidth}{>{\RaggedRight\arraybackslash}p{3.0cm}YYY}
\rowcolor{milcVino}\color{white}\bfseries Criterio & \color{white}\bfseries Lo logré & \color{white}\bfseries En proceso & \color{white}\bfseries Necesito apoyo\\
Comprensión del tema & Explico ideas centrales con mis palabras. & Reconozco ideas pero falta precisión. & Necesito repasar conceptos base.\\
Pensamiento computacional & Apliqué los 4 pilares al diseñar mi producto. & Apliqué algunos pilares. & Necesito releer la fase 2.\\
Aplicación y producto & Mi producto cumple los criterios y se puede revisar. & Producto funcional con ajustes pendientes. & Aún en construcción.\\
Reflexión 5 dimensiones & Respondí cada dimensión con honestidad. & Respondí algunas con detalle. & Mis respuestas son muy generales.\\
Triángulo de pensamiento & Conecté las 3 voces con mi trabajo real. & Conecté al menos una. & No relaciono las voces con mi proceso.\\
\end{tabularx}
\normalsize

\textbf{Hoy entendí que:}\\[1mm]
\linead{\linewidth}\\[1mm]
\linead{\linewidth}

\textbf{Mi compromiso para la próxima sesión es:}\\[1mm]
\linead{\linewidth}\\[1mm]
\linead{\linewidth}

\begin{softbox}{milcMostaza}{milcGris}{Cita de despedida · Marco Aurelio}
\textit{``El obstáculo en el camino se vuelve el camino. Lo que estorba al camino se vuelve camino.''}\\[1mm]
Cada error de hoy, bien observado, se convierte en pieza del camino que pisarás mañana.
\end{softbox}

\small
\begin{tabularx}{\textwidth}{>{\bfseries}p{3.6cm}Y}
\rowcolor{milcGradeDark!12}Nombre completo: & \linead{10cm}\\
Fecha: & \linead{4cm} \quad Curso/grupo: \linead{4cm}\\
Mi firma: & \linead{9cm}\\
\end{tabularx}
\normalsize

\begin{infoband}{milcGradeDark}
\textbf{Para la próxima sesión.} Cierra el periodo 2 (informes y comunicación profesional con IA) y abre el periodo 3 (ofimática con IA, contabilidad y emprendimiento). La sesión 1 del periodo 3 trabajará \textbf{ofimática con IA como copiloto del oficio digital cotidiano}: Google Sheets, IA aplicada a tareas reales de vida laboral próxima. Es el puente al grado 11.
\end{infoband}

\end{document}
